\documentclass[11pt]{article}
\usepackage[margin=0.75in]{geometry}
\usepackage{array}
\usepackage{enumitem}
\usepackage{xcolor}
\usepackage{amssymb}
\usepackage{hyperref}
\usepackage{listings}

\hypersetup{colorlinks=true, linkcolor=blue, urlcolor=blue}
\setlength{\parindent}{0pt}
\setlength{\parskip}{6pt}
\setlist[itemize]{topsep=3pt,itemsep=2pt,leftmargin=*}
\setlist[enumerate]{topsep=3pt,itemsep=4pt,leftmargin=*}

\lstset{
  basicstyle=\ttfamily\small,
  breaklines=true,
  frame=single,
  columns=fullflexible,
  keepspaces=true,
  showstringspaces=false,
  backgroundcolor=\color{gray!6}
}

\definecolor{predictcolor}{RGB}{255,242,200}
\definecolor{discusscolor}{RGB}{214,234,255}
\definecolor{funfactcolor}{RGB}{222,247,222}
\definecolor{debatecolor}{RGB}{255,221,221}
\definecolor{checkincolor}{RGB}{232,222,250}

\newcommand{\calloutbox}[3]{%
  \vspace{4pt}\noindent
  \begin{minipage}{\linewidth}
  \colorbox{#1}{\parbox{0.96\linewidth}{\textbf{#2}}}\\[2pt]
  \noindent\fbox{\begin{minipage}{0.94\linewidth}\vspace{2pt}#3\vspace{2pt}\end{minipage}}
  \end{minipage}
  \vspace{6pt}
}
\newcommand{\predictbox}[1]{\calloutbox{predictcolor}{PREDICTION TIME}{#1}}
\newcommand{\discussbox}[1]{\calloutbox{discusscolor}{HUDDLE UP}{#1}}
\newcommand{\funfact}[1]{\calloutbox{funfactcolor}{FUN FACT}{#1}}
\newcommand{\debatebox}[1]{\calloutbox{debatecolor}{DEBATE CORNER}{#1}}
\newcommand{\claimbox}[1]{\calloutbox{funfactcolor}{CLAIM, EVIDENCE, CONTEXT}{#1}}
\newcommand{\checkinbox}[1]{%
  \calloutbox{checkincolor}{CHECK IN: Call over an instructor}{%
    \begin{minipage}[c]{0.72\linewidth}#1\end{minipage}\hfill
    \begin{minipage}[c]{0.2\linewidth}\centering
      \fbox{\rule{0pt}{0.45in}\hspace{0.9in}}\\[2pt]
      {\small Instructor initials}
    \end{minipage}}}

\newcommand{\writebox}[1]{
  \vspace{3pt}
  \noindent\fbox{
    \begin{minipage}{0.96\linewidth}
      \textbf{#1}\\[12pt]
      \rule{\linewidth}{0.4pt}\\[14pt]
      \rule{\linewidth}{0.4pt}\\[14pt]
      \rule{\linewidth}{0.4pt}\\[14pt]
    \end{minipage}
  }
  \vspace{6pt}
}

\newcommand{\shortwritebox}[1]{
  \vspace{3pt}
  \noindent\fbox{
    \begin{minipage}{0.96\linewidth}
      \textbf{#1}\\[12pt]
      \rule{\linewidth}{0.4pt}\\[14pt]
    \end{minipage}
  }
  \vspace{6pt}
}

\begin{document}

\begin{center}
  {\LARGE MA 213 Week 9: Lab 6}\\[3pt]
  {\large Confidence Intervals and Hypothesis Tests}
\end{center}

\begin{enumerate}
  \item \textbf{Your name:} \underline{\hspace{0.3\textwidth}} \hfill \textbf{BU ID:} \underline{\hspace{0.25\textwidth}}
  \item \textbf{Partner's name:} \underline{\hspace{0.3\textwidth}} \hfill \textbf{BU ID:} \underline{\hspace{0.25\textwidth}}
\end{enumerate}

\textbf{Big idea:} Confidence intervals and hypothesis tests answer different but related questions. In this lab, you will simulate confidence interval coverage and manually build a chi-square test of independence.
\textbf{Do not use GenAI tools for this lab.} The point is to practice your own inference reasoning, coding skills, and interpretation.

\textbf{Files used:} \texttt{Titanic.csv} (the same file from Lab 3). You will also generate one population in R.

\textbf{Skills practiced:} \texttt{sample()}, \texttt{pnorm()}, \texttt{qnorm()}, \texttt{pchisq()}, \texttt{read.csv()}, \texttt{table()}, and \texttt{addmargins()}.

\textbf{Your mission:} Explain what 95\% confidence means through simulation, then test whether two categorical variables appear independent.

\section*{Icebreaker}

Before opening R: introduce yourself to your partner. Then answer: Which sounds more convincing to you, ``we estimate the proportion is between 60\% and 64\%'' or ``we reject the null hypothesis''? Why?

\shortwritebox{My answer and my partner's answer:}

\section*{Getting Started: One Proportion}

Suppose a survey asks Boston residents whether they own a cell phone. Out of 200 randomly selected residents, 127 say yes.

\begin{lstlisting}[language=R]
x <- 127
n <- 200
p_hat <- x / n

p_hat
\end{lstlisting}

\section*{Question 1. Construct and interpret a confidence interval}

Construct a 95\% confidence interval for the population proportion of Boston residents who own cell phones.

\[
\hat{p} \pm z^* \sqrt{\frac{\hat{p}(1-\hat{p})}{n}}
\]

\begin{lstlisting}[language=R]
alpha <- 0.05
# Fill in the blank with pnorm or qnorm.
z_star <- ____(1 - alpha / 2)

SE <- sqrt(p_hat * (1 - p_hat) / n)
CI <- p_hat + c(-1, 1) * z_star * SE

z_star
SE
CI
\end{lstlisting}

\writebox{State and interpret the confidence interval in context. Avoid saying there is a 95\% probability that this one interval contains the true proportion.}

\section*{Question 2. Test a hypothesis using the same sample}

Now test whether the proportion of Boston residents who own cell phones is different from 50\%.

\[
H_0: p = 0.50
\qquad
H_a: p \neq 0.50
\]

\begin{lstlisting}[language=R]
p0 <- 0.50

SE_null <- sqrt(p0 * (1 - p0) / n)
z_stat <- (p_hat - p0) / SE_null
# Fill in the blank with pnorm or qnorm.
p_value <- 2 * (1 - ____(abs(z_stat)))

z_stat
p_value
\end{lstlisting}

\writebox{What is the p-value? At $\alpha=0.05$, do you reject $H_0$? Write your conclusion in context.}

\checkinbox{\textbf{Check-in 1.} Show your confidence interval and your hypothesis test conclusion. Keep working on Question 3 until the instructor gets to you.}

\section*{Question 3. Simulate confidence interval coverage}

Let's use simulation to study whether 95\% confidence intervals for a population proportion actually contain the true proportion 95\% of the time. 

The code below creates a hypothetical population where the true proportion is known. Then it repeatedly samples from that population and builds a confidence interval each time. The ``coverage rate'' is the proportion of these intervals that contain the true proportion.

\begin{lstlisting}[language=R]
set.seed(213)

N <- 675000
pop <- c(rep("Cellphone", N * 0.9), rep("No Cellphone", N * 0.1))
true_prop <- mean(pop == "Cellphone")

true_prop
\end{lstlisting}

Run the simulation below. It builds a table with one row per simulated sample.

\begin{lstlisting}[language=R]
alpha <- 0.05
n <- 200
K <- 1000

prop_est <- rep(0, K)
lower <- rep(0, K)
upper <- rep(0, K)
included <- rep(FALSE, K)

for (k in 1:K) {
  sampled_data <- sample(pop, size = n)
  p_hat <- mean(sampled_data == "Cellphone")
  SE <- sqrt(p_hat * (1 - p_hat) / n)
  CI <- p_hat + c(-1, 1) * qnorm(1 - alpha / 2) * SE

  prop_est[k] <- p_hat
  lower[k] <- CI[1]
  upper[k] <- CI[2]
  included[k] <- CI[1] <= true_prop & true_prop <= CI[2]
}

df_table <- data.frame(
  prop_est = prop_est,
  lower = lower,
  upper = upper,
  included = included
)
\end{lstlisting}

Fill in the blanks to compute the coverage rate and the typical interval width.

\begin{lstlisting}[language=R]
# Coverage rate: the proportion of intervals that contain true_prop.
coverage_rate <- mean(df_table$____)
coverage_rate

# Typical interval width: the average width of the confidence intervals.
widths <- df_table$____ - df_table$____
typical_width <- mean(widths)
typical_width
\end{lstlisting}

Record your results in the 95\% row of the table. Then change \texttt{alpha} to 0.01 for a 99\% confidence interval, rerun the simulation, and fill in the 99\% row.

\begin{center}
\begin{tabular}{|p{0.25\textwidth}|p{0.28\textwidth}|p{0.37\textwidth}|}
\hline
\textbf{Confidence level} & \textbf{Coverage rate} & \textbf{Typical interval width}\\
\hline
95\% & & \\[20pt]
\hline
99\% & & \\[20pt]
\hline
\end{tabular}
\end{center}

\checkinbox{\textbf{Check-in 2.} Show your coverage table and explain what ``95\% confidence'' means. Keep working on Question 4 until the instructor gets to you.}

\section*{Question 4. Build a two-way table}

For the chi-square test, use the Titanic data from Lab 3. Each row of \texttt{Titanic.csv} is one passenger. The code below reads the file and cross-tabulates survival by passenger class.

\begin{lstlisting}[language=R]
titanic <- read.csv("Titanic.csv", row.names = 1)

tbl <- table(titanic$Survived, titanic$Class)
tbl
addmargins(tbl)
\end{lstlisting}

\writebox{What would independence mean in this setting? In plain language, what would it mean if passenger class and survival were independent?}

\section*{Question 5. Calculate expected counts}

For each cell in a two-way table,
Expected count = (row total)(column total) / table total.

The chi-square test requires every expected count to be at least 5.

\begin{lstlisting}[language=R]
row_totals <- rowSums(tbl)
col_totals <- colSums(tbl)
total <- sum(tbl)

expected_tbl <- outer(row_totals, col_totals) / total
expected_tbl

all(expected_tbl >= ____)
\end{lstlisting}

\textbf{Do the expected counts meet the chi-square condition?} \qquad $\square$ Yes \qquad $\square$ No

\section*{Question 6. Compute the chi-square statistic and p-value}

Compute the chi-square statistic manually, then use the chi-square distribution to find the p-value.

\[
X^2 = \sum \frac{(O-E)^2}{E}
\]

\begin{lstlisting}[language=R]
chi_stat <- sum((tbl - expected_tbl)^2 / expected_tbl)
df <- (nrow(tbl) - 1) * (ncol(tbl) - 1)
p_value <- 1 - pchisq(chi_stat, df = df)

chi_stat
df
p_value
\end{lstlisting}

If R prints a p-value of 0, the p-value is too small to display. Report it as $p < 0.001$.

\writebox{What is the p-value? At $\alpha=0.05$, do you reject $H_0$? Write your conclusion in context.}

\newpage
\section*{Question 7. Write a reusable chi-square function}

Write a function that takes a two-way table and returns the expected counts, chi-square statistic, degrees of freedom, and p-value. Check it on the class and survival table, then use it to test whether survival is independent of sex.

\begin{lstlisting}[language=R]
manual_chisq <- function(tbl) {
  # Expected counts.

  # Chi-square statistic.

  # Degrees of freedom.

  # P-value.

  list(expected = expected_tbl, chi_stat = chi_stat, df = df, p_value = p_value)
}

# Check: the results should match your answers to Questions 5-6.
manual_chisq(tbl)

# Use your function on sex and survival.
tbl2 <- table(titanic$Survived, titanic$Sex)
tbl2
manual_chisq(tbl2)
\end{lstlisting}

\writebox{Does survival appear independent of sex in the Titanic data? Support your answer with a p-value and context.}

\checkinbox{\textbf{Check-in 3.} Show your function and your conclusion about sex and survival. Turn in this worksheet with as much as you completed.}

\end{document}
